% Generated from project outputs. Do not edit by hand.
\begin{table}[!htbp]
\centering
\begin{threeparttable}
\caption{Including in-kind by subsample: 2021-2024}
\label{tab:descriptive-2021-2024-subsamples-y2}
\small
\setlength{\tabcolsep}{4pt}
\renewcommand{\arraystretch}{1.15}
\begin{tabularx}{\linewidth}{@{}Xrrrrr@{}}
\toprule
Sample & N & Mean t & SD t & Mean t+1 & SD t+1 \\
\midrule
General & 276 & 2,873.9 & 1,499.7 & 3,193.6 & 1,566.2 \\
Men & 276 & 3,246.7 & 1,632.8 & 3,638.4 & 1,675.5 \\
Women & 276 & 2,281.5 & 1,529.0 & 2,541.8 & 1,618.1 \\
Urban & 276 & 3,093.9 & 1,558.5 & 3,394.6 & 1,606.3 \\
Rural & 276 & 2,547.8 & 1,633.7 & 2,844.0 & 1,746.4 \\
Indigenous & 276 & 2,534.7 & 1,772.8 & 2,989.4 & 1,825.0 \\
Non-Indigenous & 184 & 3,031.4 & 1,560.4 & 3,230.2 & 1,512.4 \\
Bottom 99 percent & 276 & 2,695.7 & 1,282.7 & 2,993.7 & 1,307.9 \\
Bottom 90 percent & 276 & 2,231.0 & 882.1 & 2,464.9 & 869.9 \\
Bottom 50 percent & 276 & 1,259.5 & 478.0 & 1,368.4 & 448.2 \\
\bottomrule
\end{tabularx}
\begin{tablenotes}[flushleft]
\footnotesize
\item Monthly income in quetzales. Unweighted descriptive statistics of survey-weighted cohort means. N counts cohort-transition rows, not individuals or independent cohorts. SD is dispersion across cohort means, not the standard error of a mean. Both incomes must be observed. Subsamples overlap; bottom-income groups are cumulative.
\end{tablenotes}
\end{threeparttable}
\end{table}
